% Vietnamese translation for OI Public Library
% Source-File: IOI中国国家候选队论文/2003/许智磊/许智磊.doc
\documentclass[11pt]{article}
\usepackage{oipl-vn}
\usepackage{tikz}

\title{Bàn sơ về ứng dụng tư tưởng chuyển sang phần bù trong các bài toán thống kê}
\author{Xu Zhilei}
\date{}

\begin{document}
\maketitle

\origtitle{A discussion of applying complement transformation to counting problems}
\sourcepath{IOI China national candidate team papers / 2003 / Xu Zhilei / Xu Zhilei.doc}

\renewcommand{\contentsname}{Mục lục}
\renewcommand{\figurename}{Hình}
\tableofcontents

\section*{Lời nói đầu}
\addcontentsline{toc}{section}{Lời nói đầu}

Dù trong đời sống thực tế hay trong thi đấu tin học, ta thường gặp các
bài toán thống kê, tức là các bài toán đếm những đối tượng thỏa mãn một
số tính chất nào đó.

Để giải bài toán thống kê có một số phương pháp thường dùng, chẳng hạn
rời rạc hóa, cực đại hóa, chia đôi, bảng sự kiện, v.v. Dựa vào chúng, về
cơ bản ta có thể xử lý phần lớn các bài toán thống kê. Ta tạm gọi những
phương pháp ấy là các phương pháp thống kê thông thường.

Tuy nhiên, mọi sự vật trên đời vừa có tính phổ biến, vừa có tính đặc
thù. Thực sự tồn tại một bộ phận bài toán thống kê mà phương pháp thông
thường rất khó, thậm chí hoàn toàn không thể giải quyết. Với những bài
toán như vậy, cần phân tích theo tình huống cụ thể và dùng cách giải
không thông thường.

Bài viết này thảo luận một trong những cách giải không thông thường ấy:
dùng tư tưởng chuyển sang phần bù để giải bài toán thống kê.

\paragraph{Từ khóa.}
Bài toán thống kê; phương pháp thống kê thông thường/không thông thường;
liệt kê; độ phức tạp thời gian, tư duy và lập trình; chuyển sang phần bù;
đếm tổ hợp; sắp thứ tự; đối tượng liệt kê; lượng liệt kê.

\section*{Tóm tắt}
\addcontentsline{toc}{section}{Tóm tắt}

Thông qua hai ví dụ, bài toán tam giác đơn sắc (POI9714) và trò chơi hải
chiến (phỏng theo Ural 1212), bài viết khảo sát ứng dụng của tư tưởng
chuyển sang phần bù trong các bài toán thống kê, đồng thời phân tích và
so sánh tác dụng cũng như giá trị của tư tưởng này trong hai ví dụ.

Kết luận cuối cùng là: khi áp dụng vào bài toán thống kê, tư tưởng chuyển
sang phần bù thường đem lại hiệu quả rất tốt. Ta nên chú ý rèn luyện tư
duy ngược và nắm vững phương pháp thống kê không thông thường này.

\section{Mở đầu}

Ta xem bài toán thống kê là bài toán đếm những đối tượng thỏa mãn một số
tính chất nào đó. Đã là đếm thì cách giải không thể tránh khỏi việc, ít
nhiều, được xây dựng trên liệt kê. Trong rất nhiều tình huống, ``liệt
kê'' thường đồng nghĩa với hiệu quả thấp. Vì thế, trong các bài toán
thống kê thường thấy, cách giải tốt nhất đã có độ phức tạp thời gian cỡ
\(O(n^2)\) hoặc \(O(n^3)\); cũng vì vậy, quy mô dữ liệu của bài toán
thống kê thường không lớn, nói chung đều trong phạm vi \(1000\).

Nếu chỉ có vậy thì còn đỡ. Tệ hơn, với nhiều bài toán thống kê, thuật
toán thiết kế theo ý tưởng trực quan thường có độ phức tạp thời gian rất
đáng sợ: hoặc số mũ của \(n\) quá cao, hoặc bản thân \(n\) trong đề đã
rất lớn, hoặc độ phức tạp lại là biểu thức theo một đại lượng lớn khác
\(M\) (thường là kích thước bảng, số cạnh của đồ thị, v.v.). Tóm lại,
khi đó bài toán thống kê thật sự không dễ chịu.

Vì thế ta lại cần dùng nhiều kỹ thuật để giảm độ phức tạp. Các tư tưởng
rời rạc hóa và cực đại hóa, phương pháp chia đôi, bảng sự kiện, v.v. đều
rất hữu dụng. Nhờ chúng, ta đã giải thành công rất nhiều bài toán thống
kê hóc búa, đến mức các kỹ thuật ấy cũng trở thành những phương pháp
``thông thường'' hiển nhiên, và nhiều bài toán vốn khó xử cũng trở thành
bài thông thường. Nhưng rõ ràng những phương pháp này cũng có lúc không
mấy hiệu quả; khi đó ta cần một số phương pháp không thông thường để xử
lý các bài toán còn gai góc hơn.

Một trong những phương pháp không thông thường ấy là dùng tư tưởng chuyển
sang phần bù để hỗ trợ giải bài toán thống kê. Tư tưởng chuyển sang phần
bù có ứng dụng rộng rãi ở nhiều mặt; hãy xem nó có những biểu hiện thú vị
nào khi giải bài toán thống kê.

\section{Ví dụ 1: bài toán tam giác đơn sắc (POI9714 TRO)}

\subsection{Tóm tắt đề}

Trong không gian có \(n\) điểm, không có ba điểm nào thẳng hàng. Giữa hai
điểm bất kỳ đều nối bằng một đoạn thẳng đỏ hoặc đen (chỉ một đoạn, không
đỏ thì đen). Nếu ba cạnh của một tam giác có cùng màu, ta gọi tam giác ấy
là tam giác đơn sắc. Cho danh sách các đoạn thẳng màu đỏ, hãy tìm số tam
giác đơn sắc. Chẳng hạn, ở hình 1 có \(5\) điểm, \(10\) cạnh và tạo thành
\(3\) tam giác đơn sắc.

Dữ liệu vào gồm số điểm \(n\), số cạnh đỏ \(m\), và nhãn hai đỉnh của mỗi
trong \(m\) cạnh đỏ đó. Cần xuất số tam giác đơn sắc \(R\). Ràng buộc là
\[
  3 \le n \le 1000,\qquad 0 \le m \le 250000.
\]

\subsection{Phân tích ban đầu}

Một cách rất tự nhiên, ta nghĩ đến thuật toán sau: dùng một mảng
\(1000\times 1000\) để ghi màu của cạnh giữa mọi cặp điểm, rồi liệt kê
tất cả các tam giác (bằng cách liệt kê ba đỉnh), kiểm tra ba cạnh của nó
có cùng màu hay không; nếu cùng màu thì cộng vào tổng \(R\) (dĩ nhiên ban
đầu \(R=0\)).

Thuật toán này ra sao? Tạm chưa nói đến việc nó cần một mảng \(1M\) rất
xa xỉ (ở POI 97 vẫn chưa có bộ nhớ lớn), độ phức tạp thời gian của nó đã
lên tới \(O(C(n,3))\), tức là cấp \(O(n^3)\). Với bài này, khi \(n\) lớn
nhất đạt \(1000\), thật khó có thể nói đó là một thuật toán tốt.

Các kỹ thuật thông thường có dùng được cho bài này không? Rời rạc hóa và
cực đại hóa xem ra hoàn toàn không liên quan. Vì điều kiện ràng buộc của
bài rất ít, và đây cũng không phải bài đếm theo kiểu tìm giá trị tối ưu,
nên chia đôi và bảng sự kiện có vẻ cũng không dùng được.

Xem ra muốn giải bài này theo khuôn mẫu là không thể; ta cần một cách
nghĩ hoàn toàn mới.

\subsection{Suy nghĩ sâu hơn}

Phân tích thêm một chút sẽ thấy số tam giác đơn sắc trong bài này có thể
rất nhiều, nên mọi phương pháp cần liệt kê từng tam giác đơn sắc đều
không thể hiệu quả.

Liệt kê thuần túy không được, vậy đếm tổ hợp có khả thi không? Đếm tổ hợp
từ toàn cục rất khó nghĩ ra, nên ta thử liệt kê từng điểm và tìm một công
thức tổ hợp để tính số tam giác đơn sắc nhận điểm đó làm đỉnh.

Cách này dường như đã chạm đến bản chất vấn đề, vì tính bằng công thức tổ
hợp rất hiệu quả. Nhưng phân tích kỹ sẽ thấy công thức tổ hợp này rất khó
tìm. Với điểm đã liệt kê \(A\), một tam giác đơn sắc \(ABC\) có \(A\) làm
một đỉnh không chỉ cần cạnh \(AB\) và \(AC\) cùng màu, mà cạnh \(BC\) cũng
phải cùng màu với \(AB\) và \(AC\). Vì thế không thể chỉ bằng cách liệt kê
một đỉnh \(A\) mà xác định được tam giác đơn sắc.

Qua phân tích trên, ta rút ra rằng liệt kê kết hợp đếm tổ hợp có thể là
cách giải đúng, nhưng ta gặp trở ngại khi xây dựng công thức tổ hợp. Bản
chất của trở ngại này là: hai cạnh cùng màu \(AB,AC\) xuất phát từ một
đỉnh \(A\) không xác định được một tam giác đơn sắc \(ABC\), vì cạnh
\(BC\) có thể khác màu.

Nói cách khác, ta không thể thiết lập một quan hệ tương ứng xác định giữa
một cặp cạnh cùng xuất phát từ một đỉnh và toàn bộ các tam giác đơn sắc.

\subsection{Chuyển sang phần bù}

Hãy đổi góc nhìn và xét từ phía ngược lại. Vì mỗi cặp điểm đều có cạnh
nối, nên mỗi bộ ba điểm đều tạo thành một tam giác (đơn sắc hoặc không
đơn sắc). Do đó tổng số tam giác là
\[
  S=C(n,3)=\frac{n(n-1)(n-2)}{6}.
\]
Lại vì số tam giác đơn sắc \(R\) cộng với số tam giác không đơn sắc \(T\)
bằng \(S\), nên nếu tính được \(T\) thì rõ ràng \(R=S-T\). Bài toán ban
đầu vì thế tương đương với: làm thế nào tính \(T\) hiệu quả.

Thuật toán liệt kê thuần túy bị loại bỏ ngay. Vậy thuật toán liệt kê kết
hợp đếm tổ hợp, vốn thất bại trong phân tích trên, thì sao? Trở ngại ban
đầu của nó là không thể thiết lập quan hệ tương ứng xác định giữa ``một
cặp cạnh'' và ``tam giác đơn sắc''. Thế giữa một cặp cạnh có chung đỉnh
và tam giác không đơn sắc có quan hệ xác định nào không? Có. Quan hệ này
rất rõ ràng.

Ba cạnh của một tam giác không đơn sắc có cả hai màu đỏ và đen, nghĩa là
chỉ có thể có hai cạnh cùng màu và cạnh còn lại khác màu. Giả sử hai cạnh
cùng màu có đỉnh chung là \(A\), hai đỉnh còn lại là \(B\) và \(C\). Khi
đó từ điểm \(B\) nhất định xuất phát hai cạnh khác màu \(BA\) và \(BC\);
tương tự, từ điểm \(C\) xuất phát hai cạnh khác màu \(CA\) và \(CB\).

Như vậy, một tam giác không đơn sắc tương ứng với hai cặp ``cạnh khác màu
có chung đỉnh''. Mặt khác, nếu từ một đỉnh \(B\) xuất phát hai cạnh khác
màu \(BA,BC\), thì bất kể cạnh \(AC\) có màu nào, tam giác \(ABC\) chỉ có
thể là một tam giác không đơn sắc. Nói cách khác, một cặp ``cạnh khác màu
có chung đỉnh'' tương ứng với một tam giác không đơn sắc.

Rõ ràng số tam giác không đơn sắc \(T\) mà ta cần bằng một nửa tổng số
cặp ``cạnh khác màu có chung đỉnh'', ký hiệu là \(Q\). Tổng số cặp này
rất dễ tính: mỗi đỉnh có \(n-1\) cạnh, và từ dữ liệu vào ta biết số cạnh
đỏ kề với mỗi đỉnh \(i\) là \(E[i]\), nên số cạnh đen của nó là
\(n-1-E[i]\). Liệt kê đỉnh \(i\), theo nguyên lý nhân, số cặp cạnh khác
màu có \(i\) làm đỉnh chung là \(E[i](n-1-E[i])\). Do đó
\[
  Q=\sum_{i=1}^{n} E[i](n-1-E[i]).
\]

Sau khi tính được \(Q\), đáp án là
\[
  R=S-T=\frac{n(n-1)(n-2)}{6}-\frac{Q}{2}.
\]
Thuật toán này chỉ có độ phức tạp thời gian \(O(m+n)\) và độ phức tạp
không gian \(O(n)\), rất tốt.

\subsection{Tiểu kết}

Trong ví dụ này, ta đã trải qua quá trình tư duy sau:
\[
\begin{array}{c}
\text{Thuật toán liệt kê trực quan không khả thi, các kỹ thuật quen thuộc không có đất dụng võ}\\
\Downarrow\\
\text{Xét liệt kê kết hợp đếm tổ hợp}\\
\Downarrow\\
\text{Phân tích theo chiều thuận: công thức tổ hợp rất khó xây dựng}\\
\Downarrow\\
\text{Suy nghĩ từ chiều ngược lại, dùng chuyển sang phần bù và tìm được điểm đột phá.}
\end{array}
\]

Tư tưởng chuyển sang phần bù đóng vai trò then chốt ở đây: nhờ chuyển
sang phần bù, ta mới có thể thiết lập quan hệ xác định giữa ``cạnh'' và
``tam giác'', vốn không liên hệ được trong bài toán ban đầu, rồi từ đó
xây dựng công thức đếm tổ hợp. Nhờ vậy, thuật toán được chuyển từ liệt kê
thuần túy sang liệt kê kết hợp đếm tổ hợp, làm giảm mạnh độ phức tạp thời
gian và không gian. Ở đây, tác dụng của tư tưởng chuyển sang phần bù là
tạo điều kiện để tìm một thuật toán khác về bản chất.

\section{Ví dụ 2: trò chơi hải chiến (phỏng theo Ural 1212 Sea Battle)}

\subsection{Tóm tắt đề}

``Trò chơi hải chiến'' đặt các ``chiến hạm'' trên một bàn cờ ô vuông
\(N\) hàng, \(M\) cột. Một chiến hạm là một khối ô liên thông hình chữ
nhật \(X\) hàng, \(Y\) cột; mỗi ô của nó bằng đúng một ô trên bàn cờ, nên
có thể đặt chiến hạm lên bàn sao cho từng ô của chiến hạm trùng với một
ô của bàn. Khi đặt phải tuân thủ quy tắc sau: với hai chiến hạm bất kỳ,
không được có hai ô nào của chúng trùng nhau hoặc kề nhau (kề theo tám
hướng: trên, dưới, trái, phải, trên trái, trên phải, dưới phải, dưới
trái).

Hiện đã đặt \(L\) chiến hạm, thỏa mãn quy tắc đặt. Bước tiếp theo muốn
đặt thêm một chiến hạm \(P\) hàng, \(Q\) cột. Hãy tìm số cách đặt khác
nhau. Dữ liệu vào gồm \(N,M,L\), thông tin của \(L\) chiến hạm đã đặt
(tọa độ hàng cột của góc trái trên và góc phải dưới), cùng kích thước
\(P,Q\) của chiến hạm cần đặt tiếp; cần xuất số phương án \(R\). Ràng
buộc là
\[
  2 \le N,M \le 30000,\qquad 0 \le L \le 30,\qquad 1 \le P,Q \le 5.
\]
Ta xem kích thước của mọi chiến hạm đã đặt đều cùng cấp với \(P\times Q\).

Chẳng hạn, trong hình 5 đã đặt một chiến hạm \(1\) hàng \(2\) cột và một
chiến hạm \(2\) hàng \(1\) cột. Nếu cần đặt thêm chiến hạm \(1\) hàng
\(2\) cột thì có hai phương án; nếu cần đặt thêm chiến hạm \(2\) hàng
\(2\) cột thì chỉ có một phương án.

\subsection{Phân tích ban đầu}

Phương pháp liệt kê từng cách đặt rồi kiểm tra có hợp lệ hay không rõ
ràng không được, vì số cách đặt cần liệt kê đã ở cấp \(O(NM)\), chưa kể
còn phải lồng thêm quá trình kiểm tra.

Thực chất, bài gốc cho một lưới, trên đó một số vùng hình chữ nhật đã bị
chiếm dụng, và yêu cầu đặt một hình chữ nhật mới vào trong lưới sao cho
không trùng cũng không kề với các ô đã bị chiếm dụng. Đây là một bài toán
thống kê điển hình về số phương án đặt trên lưới có chướng ngại. Với một
bài toán trông kinh điển như vậy, liệu phương pháp thông thường có giải
được không? Thực ra, bằng rời rạc hóa có thể thiết kế một thuật toán chấp
nhận được. Vì đây không phải nội dung thảo luận của bài viết, ta không
nghiên cứu sâu mà chỉ nêu kết luận ban đầu: thuật toán rời rạc hóa có độ
phức tạp thời gian \(O((M+N)L)\). Tuy miễn cưỡng xử lý được bài gốc, chỉ
cần quy mô dữ liệu tăng thêm một chút là chắc chắn quá thời gian. Hơn
nữa, thuật toán rời rạc hóa có suy nghĩ tương đối phức tạp và lập trình
khá phiền.

Muốn giải bài này trọn vẹn, ta cần tìm một thuật toán mới.

\subsection{Một vài công cụ}

Trước khi suy nghĩ tiếp, ta làm rõ vài vấn đề nhỏ để dùng làm công cụ cho
phần nghiên cứu sau.

Trong một hình chữ nhật \(A\) có \(X\) hàng, \(Y\) cột, có bao nhiêu cách
đặt một hình chữ nhật \(B\) có \(P\) hàng, \(Q\) cột?

\paragraph{Kết luận 1.}
Điều kiện cần và đủ để đặt được \(B\) vào \(A\) là \(X\ge P\) và
\(Y\ge Q\). Vì thế, nếu \(X<P\) hoặc \(Y<Q\) thì số phương án bằng \(0\).
Ngược lại, góc trái trên của \(B\) có thể nằm ở các hàng từ \(1\) đến
\(X-P+1\) và các cột từ \(1\) đến \(Y-Q+1\) của \(A\), nên tổng cộng có
\[
  (X-P+1)(Y-Q+1)
\]
phương án đặt.

Góc trái trên của hình chữ nhật \(A\) là \((AX1,AY1)\), góc phải dưới là
\((AX2,AY2)\). Góc trái trên của hình chữ nhật \(B\) là \((BX1,BY1)\),
góc phải dưới là \((BX2,BY2)\). Nếu tồn tại hai ô \(a\in A\), \(b\in B\)
sao cho \(a,b\) kề nhau hoặc trùng nhau, ta nói \(A\) và \(B\) ``giao
nhau''. Làm thế nào phán đoán \(A\) và \(B\) có giao nhau hay không?

Vấn đề này hơi phức tạp hơn một chút, nhưng sau khi phân tích kỹ các
trường hợp, ta được kết luận sau.

\paragraph{Kết luận 2.}
Điều kiện cần và đủ để \(A\) và \(B\) tiếp xúc/giao nhau là
\[
\begin{aligned}
  &(BX1\le AX2+1)\ \text{and}\ (BX2\ge AX1-1)\\
  &\quad\text{and}\ (BY1\le AY2+1)\ \text{and}\ (BY2\ge AY1-1).
\end{aligned}
\]

Giả sử một hình chữ nhật đã đặt \(A\) có \(X\) hàng, \(Y\) cột, còn hình
chữ nhật mới \(B\) có \(P\) hàng, \(Q\) cột. \(B\) phải đặt như thế nào
để giao với \(A\)? Từ kết luận 2, ta trực tiếp thu được điều sau.

\paragraph{Kết luận 3.}
Mọi phương án đặt \(B\) giao với \(A\) đều nằm trong một khung hình chữ
nhật có \(2P+X\) hàng và \(2Q+Y\) cột. Trong hình 6, hình chữ nhật đen ở
giữa là \(A\); bốn hình chữ nhật xanh ở bốn góc biểu diễn các ``trường
hợp biên'' của cách đặt \(B\) có thể giao với \(A\). Dĩ nhiên, nói như
vậy vẫn chưa thật chặt chẽ, vì khung chữ nhật này có thể vượt ra ngoài
biên bàn cờ, khi đó các cạnh của nó phải được điều chỉnh vào trong biên
bàn cờ. Vì thế ta sửa kết luận 3 thành: mọi phương án đặt \(B\) giao với
\(A\) đều nằm trong một khung hình chữ nhật có nhiều nhất \(2P+X\) hàng
và nhiều nhất \(2Q+Y\) cột.

\setcounter{figure}{5}
\begin{figure}[htbp]
\centering
\begin{tikzpicture}[x=0.72cm,y=0.72cm,>=stealth]
  \draw[step=1,gray!35,very thin] (0,0) grid (8,6);
  \draw[dashed,thick,gray!70] (1,0.5) rectangle (7,5.5);
  \fill[black!72] (3.2,2.2) rectangle (4.8,3.8);
  \node[white,font=\small] at (4,3) {$A$};

  \foreach \x/\y in {1.0/4.4,5.4/4.4,1.0/0.6,5.4/0.6} {
    \fill[blue!18] (\x,\y) rectangle ++(1.6,1.0);
    \draw[blue!70!black,thick] (\x,\y) rectangle ++(1.6,1.0);
    \node[blue!70!black,font=\scriptsize] at ({\x+0.8},{\y+0.5}) {$B$};
  }

  \draw[<->,gray!70] (1,5.85) -- node[above,font=\scriptsize]
    {khung các vị trí có thể giao với \(A\)} (7,5.85);
  \node[font=\scriptsize,align=center] at (4,-0.55)
    {hình chữ nhật đen đã đặt; các hình xanh là bốn trường hợp biên};
\end{tikzpicture}
\caption{Khung chứa mọi cách đặt \(B\) có thể giao hoặc kề với \(A\).}
\end{figure}

\subsection{Chuyển sang phần bù}

Theo kết luận 1, trên bàn cờ đã cho, nếu không áp đặt hạn chế thì số
phương án đặt một hình chữ nhật, ký hiệu \(S\), có thể tính bằng công
thức. Số phương án hợp lệ \(R\) cộng với số phương án vi phạm quy tắc
\(T\) bằng \(S\), nên \(R=S-T\). Bài toán được chuyển thành: làm thế nào
tính \(T\) hiệu quả. \(T\) chính là số phương án giao với ít nhất một
chiến hạm đã đặt.

Dựa trên kết luận 3, thêm một chút phán đoán rồi kết hợp với kết luận 1,
thực ra ta đã có thể trực tiếp tính số phương án giao với một hình chữ
nhật đã đặt bất kỳ. Nhưng không thể cộng các số phương án giao với từng
hình chữ nhật đã có để làm \(T\), vì một số phương án đặt sẽ đồng thời
giao với nhiều hơn một hình chữ nhật. Chẳng hạn trong hình 7, hình chữ
nhật xanh đồng thời giao với hai hình chữ nhật đen đã đặt. Theo thuật
toán trên, phương án đặt màu xanh này sẽ bị tính lặp vào \(T\), làm cho
\(T\) lớn hơn thực tế.

\begin{figure}[htbp]
\centering
\begin{tikzpicture}[x=0.78cm,y=0.78cm,>=stealth]
  \draw[step=1,gray!35,very thin] (0,0) grid (8,5);
  \fill[black!72] (1,1.6) rectangle (2,3.6);
  \fill[black!72] (5,2.2) rectangle (7,3.2);
  \node[white,font=\scriptsize] at (1.5,2.6) {$A_1$};
  \node[white,font=\scriptsize] at (6,2.7) {$A_2$};

  \fill[blue!18] (2.4,2.0) rectangle (5.3,3.0);
  \draw[blue!70!black,thick] (2.4,2.0) rectangle (5.3,3.0);
  \node[blue!70!black,font=\small] at (3.85,2.5) {$B$};

  \draw[red!65,thick,->] (2,2.6) -- (2.4,2.6);
  \draw[red!65,thick,->] (5,2.6) -- (5.3,2.6);
  \node[font=\scriptsize,align=center] at (4,0.55)
    {\(B\) giao với cả \(A_1\) và \(A_2\), nên cộng riêng từng hình đen sẽ bị lặp};
\end{tikzpicture}
\caption{Một cách đặt \(B\) giao với hai hình chữ nhật đã đặt nên cần quy tắc
thứ tự để tránh đếm lặp.}
\end{figure}

Làm thế nào loại bỏ việc đếm lặp này? Ta dùng một phương pháp thường gặp
để loại trừ đếm lặp: sắp thứ tự. Tức là tìm cách để một phương án đặt mới
chỉ được tính vào tổng \(T\) khi xử lý hình chữ nhật đã đặt có số thứ tự
nhỏ nhất trong số các hình chữ nhật giao với nó. Ví dụ trong hình 7, nếu
quy định hình chữ nhật đen bên trái có số thứ tự nhỏ hơn, thì khi xử lý
hình chữ nhật đen bên phải, phương án đặt màu xanh giao với nó không được
tính vào \(T\), vì hình chữ nhật đen bên phải không phải hình chữ nhật đã
đặt có số thứ tự nhỏ nhất giao với phương án màu xanh. Dễ chứng minh cách
đếm này không tạo ra lặp.

Để làm được điều đó, ta chỉ cần dùng thuật toán sau. Xử lý lần lượt từng
hình chữ nhật đã đặt; giả sử số thứ tự của hình chữ nhật hiện tại là
\(i\). Liệt kê từng phương án đặt giao với nó. Với mỗi phương án, lại lần
lượt liệt kê các hình chữ nhật có số thứ tự \(1,2,\ldots,i-1\), kiểm tra
chúng có giao với phương án đang xét hay không. Nếu phát hiện có giao thì
phương án này không được tính vào \(T\); ngược lại cộng \(1\) vào \(T\).

Theo kết luận 3, các phương án đặt giao với một hình chữ nhật đã đặt có
\(X\) hàng, \(Y\) cột nằm trong một khung hình chữ nhật xung quanh nó,
khung này có nhiều nhất \(2P+X\) hàng và \(2Q+Y\) cột. Lại theo kết luận
1, trong khung đó, số cách đặt một hình chữ nhật \(P\) hàng, \(Q\) cột
nhiều nhất là \((P+X+1)(Q+Y+1)\). Vì kích thước mỗi hình chữ nhật đều ở
cùng cấp với \(P\times Q\), tổng số phương án cần xử lý có quy mô
\(O(P^2Q^2L)\). Với mỗi phương án, nhiều nhất chỉ cần liệt kê \(L-1\)
hình chữ nhật đã đặt để kiểm tra giao; theo kết luận 2, kiểm tra hai hình
chữ nhật có giao hay không có độ phức tạp \(O(1)\). Vì thế xử lý mỗi
phương án có độ phức tạp \(O(L)\), và toàn bộ thuật toán chỉ có độ phức
tạp \(O(P^2Q^2L^2)\), rất tốt, vượt xa phương pháp rời rạc hóa thông
thường. Hơn nữa, một khi đã nghĩ đến chuyển sang phần bù, phần phân tích
và thiết kế thuật toán tiếp theo đều rất tự nhiên; chương trình về cơ bản
chỉ là các vòng lặp liệt kê và các phán đoán đơn giản. Có thể nói, về độ
phức tạp tư duy và độ phức tạp lập trình, thuật toán dựa trên chuyển sang
phần bù tốt hơn nhiều so với thuật toán rời rạc hóa thông thường.

\subsection{Tiểu kết}

Trong ví dụ này, tư tưởng rời rạc hóa có thể giúp ta giải bài gốc, nhưng
độ phức tạp thời gian, tư duy và lập trình đều cao. Trong khi đó, dùng tư
tưởng chuyển sang phần bù, ta có thể thiết kế một thuật toán mới vượt
trội toàn diện so với thuật toán rời rạc hóa.

Nếu xét từ chiều thuận, lượng liệt kê của bài này quá lớn, nên cách giải
thông thường là dùng kỹ thuật rời rạc hóa để giảm lượng liệt kê. Nhưng
nếu xét từ chiều ngược lại, lượng liệt kê lại rất nhỏ. Vai trò của tư
tưởng chuyển sang phần bù ở đây là giúp ta chọn đúng đối tượng liệt kê,
từ đó giảm lượng liệt kê.

\section{Tổng kết}

Hai ví dụ trong bài viết đều dùng tư tưởng chuyển sang phần bù để giải
bài toán thống kê, và về mặt tư tưởng chúng có điểm chung rõ ràng: nếu
mục tiêu \(R\) khó giải trực tiếp, còn phần bù \(T\) của nó cùng tổng
\(S\) của \(R\) và \(T\) tương đối dễ tính, thì ta có thể xét từ mặt
ngược lại, tính \(S\) và \(T\), qua đó gián tiếp tính được \(R\).

Nhưng khi thể hiện trong hai ví dụ cụ thể, tư tưởng chuyển sang phần bù
cũng có những điểm khác nhau.

Về hiệu quả tác dụng, trong ví dụ 1, tư tưởng chuyển sang phần bù hướng
dẫn ta thiết kế một thuật toán khác về bản chất so với liệt kê thuần túy:
liệt kê kết hợp đếm tổ hợp, có thể nói là mở một con đường khác. Còn
trong ví dụ 2, nó không thay đổi bản chất thuật toán mà chỉ giảm lượng
liệt kê bằng cách thay đổi đối tượng liệt kê. Từ đó có thể thấy, hình
thức ứng dụng của tư tưởng chuyển sang phần bù trong bài toán thống kê là
rất đa dạng; nó có thể hỗ trợ ta ở mọi phương diện của việc giải bài.

Về ý nghĩa và giá trị, trong ví dụ 1, tư tưởng chuyển sang phần bù dường
như là phương pháp khả thi duy nhất để giải bài; còn trong ví dụ 2, rời
rạc hóa cũng có thể giải được, nhưng thuật toán chuyển sang phần bù tự
nhiên hơn, tốt hơn và phổ quát hơn. Từ đó có thể thấy, tư tưởng chuyển
sang phần bù không chỉ dùng được cho một số bài toán thống kê không thông
thường; ngay cả với một số bài toán mà thuật toán thông thường có thể
giải, dùng tư tưởng chuyển sang phần bù có thể vẫn làm tốt hơn.

Tóm lại, tư tưởng chuyển sang phần bù có thể được ứng dụng khá rộng rãi
trong các bài toán thống kê, và là một tư tưởng không thông thường rất
đáng nắm vững khi giải bài toán thống kê.

Tư tưởng chuyển sang phần bù thể hiện một quan niệm triết học về sự đối
lập, thống nhất và chuyển hóa lẫn nhau của mâu thuẫn. Bản chất của bài
toán thống kê là đếm trên cơ sở một loạt quan hệ và hạn chế giữa các đối
tượng đã cho, tức là trên cơ sở các mâu thuẫn. Vì vậy, vận dụng linh hoạt
tư tưởng chuyển sang phần bù trong bài toán thống kê thường có thể đem
lại hiệu quả bất ngờ. Ta nên chú ý rèn luyện năng lực tư duy ngược để
dùng tốt và dùng linh hoạt tư tưởng chuyển sang phần bù.

Điều đáng chú ý là, dùng tư tưởng chuyển sang phần bù để giải bài toán
thống kê là một phương pháp thống kê không thông thường, và quan hệ giữa
nó với một số phương pháp, kỹ thuật thống kê thông thường là quan hệ biện
chứng. Tuy trong các ví dụ của bài viết này, tư tưởng chuyển sang phần bù
có nhiều điểm sáng, nhưng không thể cho rằng phương pháp thông thường
nhất định kém phương pháp không thông thường. Phần lớn bài toán thống kê
vẫn thích hợp với phương pháp thông thường. Vì thế, chỉ khi nắm vững linh
hoạt cả phương pháp thông thường lẫn không thông thường, và biết chọn
phương pháp thích hợp cho từng bài cụ thể, ta mới có thể giải quyết bài
toán thống kê một cách thành thạo.

\end{document}
